\documentclass{beamer}
\def\presentationtype{0}
\def\corsodilaurea{Ingegneria Informatica e dell'Automazione Industriale}
\def\insegnamento{Ottimizzazione}
\def\nolezione{8}
\def\titolo{Il metodo del simplesso -- Il ``tableau'' e l'operazione ``pivot''}
\def\noattivita{1}

\def\argomenti{Programmazione Lineare}
\def\obiettivilezione{Definire il ``tableau'' e l'operazione ``pivot'' per poter sviluppare metodo del simplesso primale}
\def\nucleotematico{Richiami}
\def\contenuto{Programmazione lineare}


\usepackage[italian]{babel}
%
% Hyperref info for PDF
%
\usepackage{hyperref}
\hypersetup{%
%pdfpagemode=FullScreen,
pdfpagemode=FullScreen,
pdfstartview=Fit,
pdfsubject={\argomenti},
pdfkeywords={\insegnamento,\titolo}
}

\definecolor{ecampus@engred}{RGB}{227,28,23}
\setbeamercolor{normal text}{fg=black,bg=}
\setbeamercolor{alerted text}{fg=red}
\setbeamercolor{example text}{fg=green!50!black}
\setbeamercolor{structure}{fg=ecampus@engred,bg=}
\usetheme{default}
\useinnertheme{rounded}%{circles}
\setbeamertemplate{blocks}[default]
\setbeamercolor{block title}{bg=}
\setbeamercolor{block body}{bg=}

\ifnum\presentationtype>0 {\setbeamertemplate{enumerate item}{Es. \nolezione.\presentationtype.\arabic{enumi})}} \fi

%\setbeamertemplate{frametitle continuation}[from second][(cont'd)]
\usefonttheme[stillsansseriftext,stillsansserifsmall]{serif}
\setbeamertemplate{headline}{%
\leavevmode%
   \hspace*{4.75cm}
   \scalebox{0.90}{%
    \begin{minipage}{\linewidth}%
      \tiny
 	\begin{tabular}{ll}
	  \\
% 	  \textcolor{gray}{Classe}					& \corsodilaurea\\
%	  \textcolor{gray}{Argomento:}					& \insegnamento\\
%	  \textcolor{gray}{Lezione n}:							& \nolezione \\
	  \textcolor{gray}{Titolo:}			& \titolo \\
	  \textcolor{gray}{Attivit\`a}:		& 	  \ifnum\presentationtype=0 {Lezione \nolezione} \else {Sessione di studio \nolezione.\presentationtype} \fi
	\end{tabular}%
    \end{minipage}%
   }%
}
\setbeamertemplate{footline}[text line]{%
  \begin{minipage}{\linewidth}%
    \hfill\insertpagenumber%
%    
    \begin{center}%
      \rule{\linewidth}{0.4pt}
    \end{center}
    \vspace*{-0.65cm}
    \begin{center}%
      \scalebox{0.55}{%
	\begin{minipage}{\linewidth}%
	  \begin{center}%
	    {\tiny
	    \copyright \ 2022 Gionata Massi - \color{blue}{\href{mailto:gionata.massi@savoiabenincasa.it}{gionata.massi@savoiabenincasa.it}}\\
	    \phantom{XXX}}
	  \end{center}
	\end{minipage}%
      }%
    \end{center}%
%
  \end{minipage}%
  }
\setbeamertemplate{navigation symbols}{}

%% Save up on ink for the 4-up handouts
\mode<handout>{%
  \pgfpagesuselayout{4 on 1}[a4paper, landscape, border shrink=10mm]
  \pgfpageslogicalpageoptions{1}{border code=\pgfstroke}
  \pgfpageslogicalpageoptions{2}{border code=\pgfstroke}
  \pgfpageslogicalpageoptions{3}{border code=\pgfstroke}
  \pgfpageslogicalpageoptions{4}{border code=\pgfstroke}
}

\mode<presentation>{\AtBeginSection{%
  \begin{frame}
    \frametitle{Piano della presentazione}
    \tableofcontents[currentsection]
  \end{frame}}
}

\usepackage{microtype}
\usepackage[utf8]{inputenc}
\usepackage[T1]{fontenc}
%\usepackage[osfss]{libertine}
\usepackage{palatino}
\usepackage[scaled=.77]{beramono}
\usepackage{booktabs}
% \usepackage{attachfile}

% Toglie l'enorme spazione prima di description
%\setbeamersize{description width=0pt}

\usepackage{tikz}
\usetikzlibrary{decorations,arrows,shapes,backgrounds,matrix,positioning}
\usepackage{multirow}
\usepackage{verbatim}
\usepackage{eurosym}
\usepackage{pgfpages}

\usepackage{zref-savepos}
\newcounter{restofframe}
\newsavebox{\restofframebox}
\newlength{\mylowermargin}
\setlength{\mylowermargin}{2pt}

\newenvironment{restofframe}{%
    \par%\centering
    \stepcounter{restofframe}%
    \zsavepos{restofframe-\arabic{restofframe}-begin}%
    \begin{lrbox}{\restofframebox}%
}{%
    \end{lrbox}%
    \setkeys{Gin}{keepaspectratio}%
    \raisebox{\dimexpr-\height+\ht\strutbox\relax}[0pt][0pt]{%
    \resizebox*{!}{\dimexpr\zposy{restofframe-\arabic{restofframe}-begin}sp-\zposy{restofframe-\arabic{restofframe}-end}sp-\mylowermargin\relax}%
        {\usebox{\restofframebox}}%
    }%
    \vskip0pt plus 1filll\relax
    \mbox{\zsavepos{restofframe-\arabic{restofframe}-end}}%
    \par
}

\title{\titolo}
\author{Gionata Massi}
\date{\today}

\setbeamercolor*{frametitle}{fg=black}
\setbeamercolor*{title}{fg=black}
\setbeamerfont{frametitle}{shape=\scshape,family=\rmfamily,size=\large,series=\bfseries}

\setbeamersize{}
\makeatletter
\newcommand\thefontsize[1]{{#1 The current font size is: \f@size pt\par}}
\makeatother

\tikzstyle{nicebox}=[draw=gray!100, fill=blue!10, very thick,
rounded corners, rectangle, inner sep=4pt, inner ysep=16pt]
\tikzstyle{niceboxtitle}=[draw=gray!100, fill=white, text=black,
rounded corners, very thick, rectangle]
\newcommand\nicebox[2]{
{\centering
\begin{tikzpicture}
\node [nicebox](box){
\begin{minipage}{0.95\textwidth}\centering
\begin{minipage}{0.95\textwidth}
#2
\end{minipage}\end{minipage}};
\node[niceboxtitle, right=10pt] at (box.north west)
{\small\textbf{#1}};
\end{tikzpicture}\par}
}

\tikzstyle{modelbox}=[draw=structure!100, fill=white!50, very thick,
rounded corners, rectangle, inner sep=4pt, inner ysep=16pt, text=blue!50!black]
\tikzstyle{modelboxtitle}=[draw=structure!100, fill=white!50, text=blue!50!black,
rounded corners, very thick, rectangle]
\newcommand\modelbox[2]{
{\centering
\begin{tikzpicture}
\node [modelbox](box){
\begin{minipage}{0.95\textwidth}\centering
\begin{minipage}{0.95\textwidth}
#2
\end{minipage}\end{minipage}};
\node[modelboxtitle, right=10pt] at (box.north west)
{\small\textbf{\mbox{#1}}};
\end{tikzpicture}\par}
}

% % %
% Tipi di sessioni di studio
% % %
\def\domande{Domande}
\def\approfondimenti{Approfondimenti}
\def\esercizi{Esercizi}
\def\esempi{Esempi svolti}
\def\relazioni{Relazioni}
\def\attivitapratica{Attivit\`a pratica}
% % %

\newcommand{\generatitolo}{ %
\begin{frame}[plain]
\pdfbookmark{\ifnum\presentationtype=0 {Lezione \nolezione} \else {Sessione di studio \nolezione.\presentationtype} \fi}{\ifnum\presentationtype=0 {Lezione \nolezione} \else {Sessione di studio \nolezione.\presentationtype} \fi}
\noindent\hspace*{3.715cm}
   \scalebox{0.90}{%
    \begin{minipage}{\linewidth}%
      \tiny
 	\begin{tabular}{ll}
	  \\
% 	  \textcolor{gray}{Classe}					& \corsodilaurea\\
%	  \textcolor{gray}{Argomento:}					& \insegnamento\\
%	  \textcolor{gray}{Lezione n}:							& \nolezione \\
	  \textcolor{gray}{Titolo:}			& \titolo \\
	  \textcolor{gray}{Attivit\`a}:		& 	  \ifnum\presentationtype=0 {Lezione \nolezione} \else {Sessione di studio \nolezione.\presentationtype} \fi
	\end{tabular}%
    \end{minipage}%
   }%
   
  \begin{center}
    \large{\textbf{\insegnamento}}
  \end{center}

  \vspace*{0.5cm}
  
  \begin{center}
    \huge{\titolo}
  \end{center}
  
  \begin{center}
	  \ifnum\presentationtype=0
	    {Lezione n$^{\circ}$ \nolezione}
	  \else
	    {Sessione di studio n$^{\circ}$ \nolezione .\presentationtype}
	  \fi
  \end{center}
  
  \vspace*{0.5cm}

  \begin{center}
    \large{\textsl{Gionata Massi}}\\
    \scriptsize\color{blue}{<\href{mailto:gionata.massi@savoiabenincasa.it}{gionata.massi@savoiabenincasa.it}>}
  \end{center}
  
  \newcounter{pageleft}
  \zsaveposy{pageleft}
  \vspace*{\dimexpr\zposy{pageleft}sp-1.67cm}

  \begin{minipage}{\linewidth}%
   \hfill\phantom{\insertpagenumber}%

  \begin{center}%
      \rule{\linewidth}{0.4pt}
    \end{center}
    
    \vspace*{-0.95cm}
    
    \begin{center}%
      \scalebox{0.55}{%
	\begin{minipage}{\linewidth}%
	  %\begin{center}%
	  %  \tiny
	  %  \copyright \ 2022 Gionata Massi - \color{blue}{\href{mailto:gionata.massi@savoiabenincasa.it}{gionata.massi@savoiabenincasa.it}}\\
	  %  \color{blue}{\href{mailto:gionata.massi@savoiabenincasa.it}{gionata.massi@savoiabenincasa.it}}\\
	  %  \phantom{XXX}
	  %\end{center}
	\end{minipage}%
      }%
    \end{center}%
     \phantom{XXX}
%
  \end{minipage}%

\end{frame}
}

\renewcommand{\vec}{\mathbf}
\newcommand{\matr}{\mathbf}

\uselanguage{italian}
\languagepath{italian}
\deftranslation[to=italian]{Theorem}{Teorema}
\deftranslation[to=italian]{theorem}{teorema}
\deftranslation[to=italian]{Definition}{Definizione}
\deftranslation[to=italian]{definition}{definizione}
\deftranslation[to=italian]{Corollary}{Corollario}
\deftranslation[to=italian]{corollary}{corollario}

\let\definition\relax
\let\enddefinition\relax

%\theoremstyle{example}
\newtheorem{definition}[theorem]{\translate{Definition}}

\def\fieldN{\mathbb{N}}
\def\fieldR{\mathbb{R}}
\def\fieldZ{\mathbb{Z}}

\def\matrA{\matr{A}}
\def\vecB{\vec{b}}
\def\vecC{\vec{c}}
\def\vecX{\vec{x}}

\usepackage{cancel}
\usepackage{bookmark}

% tkiz ball item
\newcommand*\circled[1]{\raisebox{2.5pt}{\tikz[baseline=(char.base)]{
            \node[circle,ball color=purple, shade, 
 color=white,inner sep=1.2pt] (char) {\tiny #1};}}}

% tkiz rounded item
\newcommand*\rounded[1]{\tikz[baseline=(char.base)]{
            \node[draw=none,ball color=purple, shade, 
 color=white, rounded corners=3.5pt, inner sep=2.5pt] (char) {\scriptsize #1};}}

 \usepackage{etoolbox}
\makeatletter
\apptocmd{\beamer@@frametitle}{\write\@auxout{\string\@writefile{frm}{\string\frametitleentry{\the\c@framenumber}{#1}{#2}}}}{}{}
\newcommand*{\frametitleentry}[3]{\@namedef{frametitleshort#1}{#2}\@namedef{frametitle#1}{#3}}
\AtEndDocument{\if@filesw\newwrite\tf@frm\immediate\openout\tf@frm\jobname.frm\relax\fi}
\@input{\jobname.frm}
\makeatother
 
%\usebackgroundtemplate{\includegraphics[width=\paperwidth]{../img/sfondo_ecampus.png}}

\usepackage{listings}
% Extension for listings package
\lstdefinelanguage{ampl}{
  % See AMPL book (1993 edition), Appendix A
  morekeywords={%
    % Table A.1
    % Excluding symbols and "not X", where X (and not) are keywords
    if, then, else, or,  exists, forall,  and, in, within, 
    not, union, diff, symdiff, inter, cross, setof, by, 
    % Skipping: less,  sum, prod, min, max, div, mod,
    % Skipping built-in functions (Tables A-2 and A-3)
    % Section A.5 -- Declarations of model entities (subject to ->
    % subject, to)
    set, param, var, arc, minimize, maximize, subject, to, node,
    % Section A.6 -- Set declarations (not repeating set or within)
    dimen, default, card, ordered, by, reversed, circular,
    % Section A.7 -- Parameter declarations (not repeating in, default)
    param, binary, integer, symbolic, check, Infinity
    % Section A.8 -- Variable declarations
    var,  coeff, cover, obj,
    % Section A.13 -- Command language
    % Skipping most commands!
    include, model, data, let, objective, drop, restore, solve,
    solution, quit, end, option, reset, update, commands,
    % A.14 -- Imported functions
    function, pipe, 
    % A.17 synonyms for previously-defined keywords
    coef, difference, dimen, dimension, intersect, intersection, s.t., 
    % 
    % must be defined somewhere ...
    for, while, repeat, break,
    % NEW language features (http://www.ampl.com/NEW/newlang.html)
    % Grouped by page
    purge, problem, redeclare, xref,
    complements, 
    table, IN, OUT, INOUT, read, write, 
    expand,
    for, commands, repeat, until, while, break, continue,
  },
  sensitive=true,
  comment=[l]{\#},
  morecomment=[s]{/*}{*/},
  string=[b]",
  morestring=[b]',
}[keywords,comments,strings]

\definecolor{comment_c}{RGB}{60,128,49}
\definecolor{rowname_c}{RGB}{128,60,49}
\definecolor{reserved_c}{RGB}{49,60,128}

\lstset{
	language=ampl,%
	commentstyle=\color{comment_c},%
	backgroundcolor=\color{blue!15},%
    basicstyle=\small\ttfamily,%
    %emph={s.t.},emphstyle={\bfseries},%
    otherkeywords={s.t.},%
    numbers=left, numberstyle=\tiny, stepnumber=1, numbersep=5pt,%
}

\definecolor{col_set}{RGB}{128,0,0}
\definecolor{col_par}{RGB}{91,0,91}
\definecolor{col_idx}{RGB}{0,91,91}
\definecolor{col_var}{RGB}{0,0,128}

\def\insieme{%1}{\rm \textcolor{col_set}{%1}}}


\begin{document}

\generatitolo

\section{Costruzione del modello: insiemi e parametri}

\begin{frame}{Stuttura dei problemi decisionali}
I modelli decisionali hanno tutti la stessa struttura:
\begin{itemize}
\item insiemi, sui quali scorrono degli indici;
\item dati numerici (di valore noto), scalari o vettori indicizzati su un insieme o matrici a due o pi\`u dimensioni (indicizzate su due o pi\`u insiemi);
\item variabili, scalari, o vettori o matrici

\item una funzione obiettivo $z$, che dipende dai dati e dalle variabili e pu\`o essere minimizzata o massimizzata;
\item un certo numero di vincoli, di uguaglianza (equazioni) o disuguaglianza
(disequazioni), singoli oppure raccolti in famiglie indicizzate su uno o pi\`u insiemi.
\end{itemize}
\end{frame}

\begin{frame}{Mix ottimo di produzione: modello basato su insiemi}
\definecolor{col_set}{RGB}{128,0,0}
\definecolor{col_par}{RGB}{91,0,91}
\definecolor{col_idx}{RGB}{0,91,91}
\definecolor{col_var}{RGB}{0,0,128}
\begin{columns}%{}

\column{.49\textwidth}

\begin{minipage}{\textwidth}
\footnotesize{
\begin{tabular}{lp{0.85\textwidth}}
\color{col_set}$P$		&Insieme dei prodotti\\
\color{col_set}$R$		&Insieme delle risorse\\
~\\
$\textcolor{col_par}{a}_{\textcolor{col_idx}{ij}}$	&Assorbimento della risorsa $i\in R$ da un prodotto $j\in P$\\
$\textcolor{col_par}{b}_{\textcolor{col_idx}{i}}$	&Disponibilit\`a della risorsa $i\in R$\\
$\textcolor{col_par}{c}_{\textcolor{col_idx}{j}}$	&Profitto unitario prodott $j\in P$\\
~\\
$\textcolor{col_var}{x}_{\textcolor{col_idx}{j}}$	&Livello di produzione prodotto \mbox{$j\in P$}\\
~\\
\end{tabular}
\begin{tabular}{ll}
\hspace*{0.5\textwidth}&\\
\textcolor{col_set}{\rule{1 em}{1 em} Insiemi} &
\textcolor{col_idx}{\rule{1 em}{1 em} Indici} \\
\textcolor{col_par}{\rule{1 em}{1 em} Parametri} &
\textcolor{col_var}{\rule{1 em}{1 em} Variabili}
\end{tabular}
}
\end{minipage}

\column{.51\textwidth}

\def\funzioneobiettivo{{\displaystyle \sum_{\textcolor{col_idx}{j} \in \textcolor{col_set}{P}}}
	\textcolor{col_par}{c}_{\textcolor{col_idx}{j}}\,\textcolor{col_var}{x}_{\textcolor{col_idx}{j}}}
\def\jinP{\textcolor{col_idx}{j} \in \textcolor{col_set}{P}}
\def\iinR{\textcolor{col_idx}{i} \in \textcolor{col_set}{R}}


\fbox{
{\footnotesize
$\begin{array}{rll}
\max z= & \funzioneobiettivo\\% Funzione obiettivo
	       & {\displaystyle \sum_{\textcolor{col_idx}{j} \in \textcolor{col_set}{P}}}
	       \textcolor{col_par}{a}_{\textcolor{col_idx}{ij}} \,
	       \textcolor{col_var}{x}_{\textcolor{col_idx}{j}} \leq \textcolor{col_par}{b}_{\textcolor{col_idx}{i}} & \iinR\\
  & \textcolor{col_var}{x}_{\textcolor{col_idx}{j}}\geq 0 &  \jinP, \, \iinR
\end{array}$
}
}
\end{columns}%{}
\end{frame}


\section{Insiemi, indici, parametri e iteratori}

\begin{frame}[allowframebreaks,fragile]{Insiemi}

\framebox{
\parbox[c][24pt]{0.85\textwidth}{
\hspace{6pt} {\tt set} {\it name} {\it alias} {\it domain} {\tt,}
{\it attrib} {\tt,} \dots {\tt,} {\it attrib} {\tt;}
}}

\begin{description}

\item[{\it name}] \`e il nome dell'insieme;

\item[{\it alias}] opzionale, altro nome dell'insieme;

\item[{\it domain}] opzionale, specifica meglio il dominio;

\item[{\it attrib}, \dots, {\it attrib}] attributi opzionali.
\end{description}

\framebreak

{\small Attributi opzionali}

\vspace*{-8pt}

{\footnotesize \begin{description}
\item[{\tt dimen} $n$]\hspace*{0pt}\\
la cardinalit\`a $n$ dell'insieme;
\item[{\tt within} {\it expression}]\hspace*{0pt}\\
un insieme che include l'insieme che si sta definendo;
\item[{\tt:=} {\it expression}]\hspace*{0pt}\\
i membri dell'insieme (enumerazione o propriet\`a);
\item[{\tt default} {\it expression}]\hspace*{0pt}\\
gli elementi dell'insieme se
non specificati nella sezione dati.
\end{description}}

\begin{block}{Insiemi ``Prodotti'' e ``Risorse''}
\begin{lstlisting}
set Prodotti;
set Risorse;
\end{lstlisting}
\end{block}
\end{frame}

\begin{frame}[allowframebreaks,fragile]{Indici}
Due forme delle espressioni di indicizzazione, che permettono di scorrere
gli indici:
$$
{
\begin{array}{l}
\mbox{{\tt\{} {\it entry}$_1${\tt,} {\it entry}$_2${\tt,} \dots{\tt,}
{\it entry}$_m$ {\tt\}}}\\
\mbox{{\tt\{} {\it entry}$_1${\tt,} {\it entry}$_2${\tt,} \dots{\tt,}
{\it entry}$_m$ {\tt:} {\it predicate} {\tt\}}}\\
\end{array}
}
$$
dove {\it entry}{$_1$}, {\it entry}{$_2$}, \dots, {\it entry}{$_m$}
sono elementi di indicizzazione e {\it predicate} \`e un'espressione
logica opzionale che deve essere verificata dagli indici.

Ogni {\it indexing entry} \`e in una delle forme:
$$
{
\begin{array}{l}
\mbox{$i$ {\tt in} $S$}\\
\mbox{{\tt(}$i_1${\tt,} $i_2${\tt,} \dots{\tt,}$i_n${\tt)} {\tt in}
$S$}\\
\mbox{$S$}\\
\end{array}
}
$$
dove $i_1$, $i_2$, \dots, $i_n$ sono indici, $S$ \`e un insieme base.

Sono usate, ad esempio, per indicizzare i parametri di un modello e di un iteratore

\begin{lstlisting}
{j in Prodotti};
{j in Prodotti : j <> prod_1};
\end{lstlisting}
\end{frame}

\begin{frame}[allowframebreaks,fragile]{Parametri}

\framebox{
\parbox[c][24pt]{0.85\textwidth}{
\hspace{6pt} {\tt param} {\it name} {\it alias} {\it domain} {\tt,}
{\it attrib} {\tt,} \dots {\tt,} {\it attrib} {\tt;}
}}

\begin{description}

\item[{\it name}] \`e il nome del parametro;

\item[{\it alias}] opzionale, altro nome del parametro;

\item[{\it domain}] opzionale, specifica meglio il dominio;

\item[{\it attrib}, \dots, {\it attrib}] attributi opzionali.
\end{description}

\framebreak

{\small Attributi opzionali}

\vspace*{-8pt}

{\footnotesize \begin{description}
\item[{\tt integer}]\hspace*{0pt}\\
il parametro \`e un numero intero;
\item[{\tt binary}]\hspace*{0pt}\\
il parametro \`e binario, 0/1;
\item[{\tt symbolic}]\hspace*{0pt}\\
il parametro \`e simbolico;
\item[{\it relation expression}]\hspace*{0pt}\\
(dove {\it relation} \`e un simbolo tra: {\tt<}, {\tt<=}, {\tt=}, {\tt==},
{\tt>=}, {\tt>}, {\tt<>}, {\tt!=})\\
una condizione che vincola il parametro o i suoi membri a soddisfare
quella condizione;
\item[{\tt in} {\it expression}]\hspace*{0pt}\\
un insieme che vincola il parametro o i suoi membri ad appartenere a tale
insieme;
\item[{\tt:=} {\it expression}]\hspace*{0pt}\\
un valore assegnato al parametro o ai sui membri;
\item[{\tt default} {\it expression}]\hspace*{0pt}\\
un valore assegnato al parametro o ai suoi membri se
non specificati nella sezione dati.
\end{description}}

Per riferirsi a una componete di un parametro definito come
matrice si usa il costrutto:
$$
\mbox{{\it name}{\tt[}$i_1${\tt,} $i_2${\tt,} \dots{\tt,} $i_n${\tt]}}
$$
dove {\it name} \`e il nome del parametro e $i_1$, $i_2$,
\dots, $i_n$ sono indici.

\begin{block}{Parametri ``profitto'', ``disponibilita'' e ``assorbimento'''}
\begin{lstlisting}
param profitto {j in Prodotti};
param disponibilita {i in Risorse};
param assorbimento {i in Risorse, j in Prodotti};
\end{lstlisting}
\end{block}
\end{frame}

\begin{frame}{Iteratori e quantificatori}
\noindent\hfil
\begin{tabular}{@{}lll@{}}
{\tt sum}&somma&$\displaystyle\sum_{(i_1,\dots,i_n)\in\Delta}
f(i_1,\dots,i_n)$\\
{\tt prod}&prodotto&$\displaystyle\prod_{(i_1,\dots,i_n)\in\Delta}
f(i_1,\dots,i_n)$\\
{\tt min}&minimo&$\displaystyle\min_{(i_1,\dots,i_n)\in\Delta}
f(i_1,\dots,i_n)$\\
{\tt max}&massimo&$\displaystyle\max_{(i_1,\dots,i_n)\in\Delta}
f(i_1,\dots,i_n)$\\
{\tt forall}&$\forall$&$\displaystyle
\forall(i_1,\dots,i_n)\in\Delta[f(i_1,\dots,i_n)],$\\
{\tt exists}&$\exists$&$\displaystyle
\exists(i_1,\dots,i_n)\in\Delta[f(i_1,\dots,i_n)],$\\
\end{tabular}

\noindent \\~\\ dove $i_1$, \dots, $i_n$ sono gli indici degli insiemi,
$\Delta$ \`e il dominio,
$f(i_1,\dots,i_n)$ \`e l'integrando, un espressione numerica 
il cui valore dipende dagli indici.
\end{frame}

\section{Descrizione di un modello}

\begin{frame}{Sezioni}
La dichiarazione del programma lineare \`e divisa nelle sezioni modello
(\emph{model section}) e dati (\emph{data section}).

\begin{description}
\item [{Model section}] descrive il modello matematico mediante
le dichiarazioni degli oggetti del modello; essi che sono comuni a
tutti i problemi definiti sullo stesso modello.
\item [{Data section}] opzionale, contiene i dati specifici di una
particolare istanza del problema.
\end{description}

Le sezioni ``modello'' e ``dati'' possono essere memorizzate in un
unico file o in due file separati.

La seconda opzione permettedi avere un numero arbitrario di sezioni
dati differenti da usare con la stessa sezione dati.
\end{frame}

\section{Implementazione del modello del MIX OTTIMO}

\begin{frame}{Formulazione GNU Mathprog del mix ottimo -- Model}
\lstinputlisting{allegati/modello_algebrico.mod}
\end{frame}

\begin{frame}{Formulazione GNU Mathprog del mix ottimo -- Data}
\lstinputlisting{allegati/dati_mix_ottimo.dat}
\end{frame}

\section{Per approfondire}

\begin{frame}{Per approfondire}

\beamertemplatebookbibitems
\begin{thebibliography}{References}
\bibitem{Cordone}R. Cordone.\newblock
\textit{Il linguaggio GNU MathProg con esempi commentati e risolti}.\newblock
2010. Disponibile su internet.

\bibitem{gmplManual} A. Makhorin.\newblock
\textit{Modeling Language GNU MathProg -- Language Reference}
\newblock 2013. Distribuito con GLPK.

\bibitem{amplbook} R. Fourer, D. M. Gay, B. W. Kernighan.%
\newblock \emph{A Modeling Language for Mathematical Programming.}%
\newblock Management Science, 36, 1990. Disponibile su internet.
\end{thebibliography}
\end{frame}

\end{document}
